\section{Scope and configured boundaries}
\label{sec:discussion}

\subsection{Configured forward-model scope}

The coherent-optical default is vectorial Debye, with explicit scalar paraxial and other fidelity knobs available through backend selectors when configured. Fresnel-reference and high-numerical-aperture dipole-collection controls are opt-in features used by the native interferometric scattering microscopy reference-check profile. The vectorial fluorescence backend includes a vectorial point-spread-function hook, directional bright-to-dark and dark-to-bright state kinetics, spectral scaling, bleaching controls, and detector-domain photon/electron/count propagation; proxy fluorescence modes remain selectable for lightweight diagnostics. Gibson--Lanni-style coverslip mismatch is exposed as a configured pupil optical-path-difference term; zero mismatch reduces to the ordinary backend and nonzero mismatch is recorded in profile metadata. Optional particle-volume helpers expose z-stack, confocal axial weighting, light-sheet axial weighting, and holotomography-style phase projection stacks. The manuscript's default comparison tables state which configured profile they used rather than treating every optional backend as implicit.

The default optical path is single wavelength. Same-scene spectral rendering is also available through caller-configured channels and an automatic broadband-quadrature channel generator. That spectral path renders one latent scene at each wavelength, integrates detector channels from ideal signal/reference frames, and applies detector noise after spectral integration. Fluorescence carries separate excitation and emission wavelengths; backend metadata distinguishes parametric point-spread-function profiles from the vectorial photophysics profile rather than collapsing both into one fidelity class.

The default ranking tables use the single-frame contract unless a sequence-aware contract is selected. For a static particle observed over $K$ independent, identically configured frames, Fisher information scales by $K$ and the Cram\'{e}r--Rao bound improves by $\sqrt{K}$. Syniscopy also exposes sequence Fisher accumulation and a dynamic Bayesian information-filter CRLB with Brownian process covariance, initial covariance, frame rate, state axes, and unit metadata. The output provenance records whether the row used \texttt{per\_frame}, \texttt{static\_cumulative}, or \texttt{dynamic\_bayesian\_information\_filter}; per-frame CRLB is therefore a selected contract, not a missing estimator layer.

The comparison claims in this paper are contract-scoped: the emitted status fields determine whether a row is validated evidence or only a diagnostic record, not scalar claims over all data-collection setups. A microscope can dominate in Contract-LP, differ in Contract-LZ, and invert under Contract-Q because each contract encodes distinct assumptions. That is the comparison design: Contract-LP, Contract-LZ, and Contract-Q answer different experimental questions and should not be collapsed into one universal ordering. The instrument configuration is part of this scope: because the comparison unit is a microscope rather than a bare modality, two configurations of the same modality can rank on opposite sides of a contract, and the configured-profile tables fix one instrument only to isolate the contrast modality.

\subsection{Validation and acquisition metadata}

The sample environment includes a structured thin-film or support layer with explicit height and material maps. The environment architecture already supports multilayer and patterned stacks, and additional layered/contamination/changing-interface support is represented through the extensible pattern and material interfaces. Sub-feature roughness is implemented via configurable interface roughness/speckle fields \texttt{sample\_environment\_pattern\_roughness\_*}, including static, flicker, and source-matched modes and roughness-source coupling options \texttt{independent}, \texttt{coherent\_amplitude}, \texttt{field\_weighted}, \texttt{scene\_weighted}, and \texttt{channel\_weighted} (default: \texttt{channel\_weighted}). Drift, vibration, empirical shading, fixed-pattern noise, calibrated gain/scale terms, hot pixels, dark current, and scanning electron microscopy scan-line noise are implemented as configurable correction channels with optional temporal realization and per-frame composition semantics. Detector quantum efficiency and EMCCD excess-noise behavior are implemented as separate public top-level parameters in the configuration schema (alongside scan-noise and fixed-pattern terms).

Transmission electron microscopy, scanning electron microscopy, and fluorescence now use explicit backend selectors rather than implicit fidelity assumptions, and the selected fidelity is recorded per row rather than assumed. Backend fidelity is a deliberate configurable axis with a computational-cost trade-off, so the software default and the paper-facing comparison choice are intentionally distinct. Transmission electron microscopy defaults to its physics-based \texttt{multislice\_physical} backend, which is also the comparison backend. Scanning electron microscopy defaults to a fast physics-based lightweight transport rather than to the full Monte Carlo path: full material-resolved \texttt{monte\_carlo\_physical} transport is computationally expensive and is opt-in, not the default, because it is unnecessary for routine exploratory comparison. The paper's headline electron-microscopy comparison rows explicitly select the highest-fidelity physical backends --- \texttt{multislice\_physical} for transmission electron microscopy and material-resolved \texttt{monte\_carlo\_physical} transport for scanning electron microscopy --- and lower-fidelity proxies and the full Monte Carlo path are both selectable, so a reader knows exactly which fidelity produced each row. The TEM implementation is physics-based but remains labelled reference-unvalidated unless an external multislice-backend validation hash is supplied. The SEM implementation is checked against published-style transport anchors: screened-Rutherford elastic scattering with Joy--Luo stopping is the current baseline, while the Browning empirical Mott surrogate is available as an explicit sensitivity option and is not promoted to the default because it worsens the present backscatter/range anchor residuals under the simplified kernel geometry. Fluorescence can route to the vectorial photophysics backend. The profile-card and model-card contracts record \texttt{backend\_fidelity\_level} and \texttt{reference\_backend\_metadata}; a row is labelled reference-validated only when the selected backend supplies the corresponding validation metadata. Electron-microscopy and optical rows are comparable only in the contract-level information sense: the same latent particle state and parameter frame are evaluated under declared measurement-domain and noise/count/dose assumptions. Vacuum, preparation, liveness, destructive acquisition, and sample-compatibility constraints therefore limit experimental-design interpretation even when they do not change the algebraic Fisher diagnostic. Appendix~\ref{sec:source-use-audit} separates source-use categories: parameter-matched validation, literature-scale checks, modality-principle citations, and dimensional-metrology references are not collapsed into a single validation claim.

Signal, information, temporal, and assignment-ambiguity support factors are bounded support scores by default. Syniscopy also exposes calibrated-probability semantics through explicit Platt-logistic or isotonic calibration parameters plus reliability metadata. The policy prevents unsupported known-object pixels from becoming ordinary background negatives; calibrated probability language is gated on supplied calibration metadata. The default ambiguity factor is the uniform-assignment special case; non-uniform ambiguity scoring uses explicit positions or masks for competing particles in the frame.

\subsection{Domain boundaries}

Syniscopy models discrete nanoscale scatterers on mounting interfaces or in suspension. Cell morphology, subcellular structure, tissue context, live-cell organelle tracking, cell-population statistics, cell-boundary segmentation, and tissue segmentation are separate application domains with different latent-state schemas. High-angle annular dark-field scanning transmission electron microscopy, atomic force microscopy, stochastic optical reconstruction microscopy, photoactivated localization microscopy, and MINFLUX are separate measurement methods. SEM transport is represented through material-resolved physical Monte Carlo transport, explicitly selected lower-fidelity raster/proxy transport, and reference-kernel interfaces; reference validation is issued only when required kernel metadata and hashes are present.
